\section{模型求解流程与实现细节}

本文的求解采用“数据校验—日内优化—月度覆盖—结果回代”的四阶段流程。首先对原始附件进行单位统一、缺失值检查和逐日聚合；其次以每一天的小时到货序列为输入，分别求解问题一和问题二；再次将问题二得到的最低需求向量传递给问题三；最后把出勤矩阵回代到班次人数中，检查每一个小时的流量守恒、截止时限和月度劳动规则。该流程的优点是每个子模型都有清晰的业务含义，任一阶段出现异常时可以定位到具体约束，而不是只观察一个最终目标值。

在实现上，班次起点的候选集合为$S=\{0,1,\ldots,16\}$。由于一个班次持续8小时，起点为16的班次覆盖16--23时，恰好不跨日。对给定的起点组合，$y_{ds}$决定静态人力投入，$p_{dh}$和$b_{dh}$负责描述动态货物流。与直接使用“总货量除以总工时”的估算相比，流量模型能够识别早班空闲、午间拥堵和晚班补偿三类不同情形。

\begin{table}[htbp]
\centering
\caption{模型求解的输入、决策与校验输出}
\label{tab:workflow}
\begin{tabular}{p{.22\textwidth}p{.28\textwidth}p{.35\textwidth}}
\toprule
阶段 & 主要输入/决策 & 输出与校验内容\\
\midrule
数据预处理 & 720条小时到货记录 & 单位、日期索引、日总量与峰值小时\\
日内优化 & $x,y,p,b,z$ & 5个班次、人数、逐小时处理量与库存\\
月度排班 & $w_{id}$ & 招工总数、每日到岗人数、休息窗口\\
结果回代 & 三类模型输出 & 货物清零、16时截止、23个工日\\
\bottomrule
\end{tabular}
\end{table}

为避免整数求解器将数值误差误判为业务可行，程序在求解完成后重新计算每个小时的容量上限，并以$10^{-6}$为连续变量误差阈值；整数变量则直接检查是否为整数。只有同时通过库存非负、日末清零、截止约束、每日覆盖和滚动窗口检查的解，才被写入结果文件。

\section{问题一的结构性分析}

问题一的一个重要性质是：在不考虑班次数量限制时，日总人力存在容量下界
\begin{equation}
 L_d^{(1)}=\left\lceil \frac{\sum_{h=0}^{23}a_{dh}}{200}\right\rceil .
\end{equation}
然而该下界通常不能直接达到。原因在于每名工人的200件能力必须分布在连续8小时内，不能把凌晨未使用的处理能力搬运到白天；同时，货物到达具有时序性，$p_{dh}\le a_{dh}+b_{d,h-1}$又限制了处理量不能提前发生。因此实际最优值由“总量下界”和“时间覆盖下界”共同决定。

从排班角度看，5个班次并不等价于5个任意的8小时容量块。若所有班次集中在高峰时段，虽然高峰小时容量充足，却可能导致低峰时段没有足够人员清空积压；若班次过于分散，又会增加交叠所造成的重复覆盖。模型通过$x_{ds}$选择起点，通过库存变量自动衡量不同起点组合的真实效果。

为了说明这一点，定义小时利用率
\begin{equation}
 u_{dh}=\frac{p_{dh}}{25\sum_s\delta_{sh}y_{ds}},
\end{equation}
当分母为零时令$u_{dh}=0$。利用率低并不必然表示排班不合理，因为该小时可能尚未有足够货物到达；真正需要重点关注的是高到货小时的利用率和日末库存。由此，本文不以“每小时满载”为目标，而以满足时序可行的最少人数为目标。

\section{问题二的服务水平解释}

问题二增加的16时约束本质上是一个带优先级的服务水平约束。把0--12时到货视为早间订单集合，要求其在前16个小时内完成处理。该要求会改变最优排班的形状：即使全天总货量不变，若早间到货集中在前几小时，模型也必须提前配置更多人力；反之，若早间到货较少，后续班次仍可承担主要处理任务。

本文采用累计处理量形式而不是简单要求$\sum_{h=0}^{15}p_{dh}=\sum_{h=0}^{12}a_{dh}$，原因是13--15时到货也可能在16时前被处理，因而截止时刻的实际处理量可以超过早间到货量。库存守恒保证了多处理的货物一定已经到达，避免了截止约束带来的“虚假提前处理”。

低效率小时的处理采用工人数分解。对同一班次而言，$z_{dsk}$允许不同工人在不同相对小时处于低效率状态，从而描述错峰休息或设备操作。若强制所有人同时低效，容量会出现不必要的同步下降，并夸大问题二的人力需求；若完全忽略低效率小时，又会高估实际处理能力。两种极端之间的整数分配正是该模型相较于粗略经验公式的改进。

\begin{table}[htbp]
\centering
\caption{问题二相对问题一的人力增量解释}
\label{tab:increment}
\begin{tabular}{lrrp{.40\textwidth}}
\toprule
指标 & 问题一 & 问题二 & 解释\\
\midrule
总人日 & 10791 & 11951 & 低效率损失与早间截止共同作用\\
平均日人数 & 359.70 & 398.37 & 约束从全天清零提升为分阶段清零\\
峰值人数 & 537 & 581 & 峰值日同时承受高货量和高服务压力\\
增幅 & -- & 10.75\% & 可作为服务水平成本的量化指标\\
\bottomrule
\end{tabular}
\end{table}

\section{问题三的可行性证明与调度含义}

问题三不是简单地把每天需求相加，而是一个带固定月度工日的集合覆盖问题。$N=581$首先由第12日峰值需求给出不可突破的下界；而$23N$与总需求之间的差额说明，固定工日制度必然产生一部分结构性冗余。该冗余不能通过减少招工消除，只能通过休息日错峰、机动岗位和班次内再分配加以吸收。

滚动8日约束具有很强的可解释性。式\eqref{eq:window}要求每个8日窗口至少出现1个休息日，因此它同时排除了8天连续工作，也避免了把23个工日集中安排在月初或月末。与逐人贪心排班相比，整体0--1模型可以把峰值日期作为覆盖约束统一处理，在第12日等紧张日期优先保证足够人员，再用相邻日期的休息日进行平衡。

在实际管理中，可将求解结果分成三类人员：高峰日固定骨干、随日期轮换的普通作业人员和承担超额覆盖的机动人员。若临时出现某日需求上调，只需重新求解覆盖约束，通常不必改动日内货物流模型；若出现人员请假，则可以把对应的$w_{id}$固定为0后重新优化，从而形成可操作的滚动重排机制。

\section{模型结果的管理解释}

从结果看，问题一给出的10791人日代表“只要求当天清零”的最低用工基线；问题二的11951人日则代表在服务水平和低效率现实条件下的计划用工基线；问题三的581名招工人数是满足月度劳动规则的最小编制。三者分别对应运营、计划和人事三个层次，不能混为同一个指标。

图表分析还揭示了一个重要现象：峰值人数并不完全由日总货量决定。某些日期总货量并非最高，但由于到货早、截止压力大，问题二的需求增量反而更明显。因此实际管理中应同时监控日总量、早间到货占比、峰值小时和16时前剩余库存，而不能只根据日总量制定班表。

敏感性实验表明，当低效率损失分别缩放为基准的0.95、1.00和1.05时，总人日分别为12576、11951和11383，峰值人数分别为611、581和553。该实验说明效率参数对人力计划具有明显的非线性影响：当有效处理能力降低时，峰值日期首先触发额外人员配置，随后才表现为一般日期的整体增加。因此设备维护、培训和错峰休息具有直接的人力节约价值。

\begin{table}[htbp]
\centering
\caption{效率参数敏感性结果及决策建议}
\label{tab:sensitivity2}
\begin{tabular}{lrrp{.42\textwidth}}
\toprule
效率缩放 & 总人日 & 峰值人数 & 管理含义\\
\midrule
0.95 & 12576 & 611 & 低效损失扩大，应预留峰值机动人员\\
1.00 & 11951 & 581 & 基准排班与月度581人编制\\
1.05 & 11383 & 553 & 通过培训和设备改善可降低用工\\
\bottomrule
\end{tabular}
\end{table}
